\chapterimage{figs/chapterhead.png}
\chapter{Shell Usage}
\label{chap:shell-usage}

This chapter documents the current terminal shell as implemented in
\texttt{source/applications/shell/}. It is based on code inspection of:
\begin{itemize}
\item \texttt{GenesysShell.cpp};
\item \texttt{GenesysShell.h};
\item \texttt{main.cpp};
\item \texttt{BaseGenesysTerminalApplication.cpp};
\item the current \texttt{genesys\_shell} CMake target.
\end{itemize}

\noindent\textbf{Objective.} Document the supported terminal workflow for users
who want to inspect a model, load a file, run a simulation, or automate a
small sequence of shell commands.

\noindent\textbf{Audience.} Readers who are comfortable with a terminal and
want a lightweight command surface instead of the GUI.

\noindent\textbf{Scope.} Cover build and launch, startup behavior, the current
command families, model loading and validation, simulation control, scripted
automation, and the current limits of the shell.

\noindent\textbf{Status.} Confirmed by code inspection and aligned with the
current shell target. The figure blocks are structural placeholders until the
final screenshots are produced.

The chapter follows this sequence:
Figure~\ref{fig:shell-session},
Figure~\ref{fig:shell-command-flow},
Figure~\ref{fig:shell-model-loaded}, and
Figure~\ref{fig:shell-terminal-result}.

\section{Build and launch}
\label{sec:shell-build-and-launch}

The shell target is named \texttt{genesys\_shell}. The dedicated preset is
\texttt{genesys\_shell}, and its build directory is configured under
\path{build/genesys_shell}.

The current build entry points are:
\begin{itemize}
\item \texttt{cmake --preset genesys\_shell};
\item \texttt{cmake --build --preset genesys\_shell --parallel "\$(nproc)"}.
\end{itemize}

The executable starts by constructing \texttt{GenesysShell}, setting the
default trace handlers, and then invoking \texttt{main(argc, argv)}. During
startup, the shell temporarily lowers the trace level, tries to auto-load
plugins from \texttt{autoloadplugins.txt}, restores the internal trace level,
and then enters the prompt loop.

If \texttt{autoloadplugins.txt} is missing or incomplete, the shell prints an
error and continues. The shell still starts; plugin loading failure does not
abort the process.

\begin{figure}[htbp]
\centering
% FIGURE-SPEC-BEGIN
% ID: fig:shell-session
% Source of truth:
%   - source/applications/shell/main.cpp
%   - source/applications/shell/GenesysShell.cpp
%   - source/applications/shell/GenesysShell.h
%   - source/applications/shell/CMakeLists.txt
% Purpose:
%   Show a complete terminal session from startup prompt through a clean exit.
% Required visual content:
%   - the shell prompt
%   - the startup message
%   - plugin autoload output or its failure message
%   - at least one successful command and the quit command
% Accessibility description:
%   A terminal screenshot that proves the shell starts, accepts commands, and exits cleanly.
% Validation:
%   The visible prompt and messages must match the current shell implementation.
% FIGURE-SPEC-END
\figureplaceholder{figs/generated/shell-session}
\caption{Complete terminal session for the Genesys shell.}
\label{fig:shell-session}
\end{figure}

Figure~\ref{fig:shell-session} is the reference view for the rest of the
chapter. It should show the prompt \texttt{\$genesys>} and a normal shutdown
path.

\section{Command surface}
\label{sec:shell-command-surface}

The shell accepts a single command word followed by optional arguments. The
first word is matched case-insensitively; the remaining words are preserved as
arguments and rejoined only by the commands that explicitly do so.

The current command families are:
\begin{itemize}
\item \texttt{help} and \texttt{exit};
\item metadata and identity: \texttt{facade}, \texttt{info}, \texttt{licence},
and \texttt{trace};
\item plugin management: \texttt{plugin};
\item model lifecycle: \texttt{model}, \texttt{load}, and \texttt{save};
\item simulation control: \texttt{sim}, \texttt{simul}, and \texttt{config};
\item inspection helpers: \texttt{data}, \texttt{component}, and
\texttt{experiment};
\item automation helpers: \texttt{execute-script} and \texttt{bash}.
\end{itemize}

The command families are intentionally conservative. The shell does not invent
new semantics on top of the underlying simulator facade; it forwards to the
current code paths and prints textual diagnostics when a request cannot be
fulfilled.

\begin{figure}[htbp]
\centering
% FIGURE-SPEC-BEGIN
% ID: fig:shell-command-flow
% Source of truth:
%   - source/applications/shell/GenesysShell.cpp
%   - source/applications/shell/GenesysShell.h
% Purpose:
%   Show the current command families and how they map to the shell handlers.
% Required visual content:
%   - the main command groups
%   - the special grouped commands with topic help
%   - the automation and inspection branches
% Accessibility description:
%   A diagram that helps the reader understand how the shell routes each command family.
% Validation:
%   The groups and labels must match the current command definitions.
% FIGURE-SPEC-END
\figureplaceholder{figs/generated/shell-command-flow}
\caption{Current Genesys shell command families.}
\label{fig:shell-command-flow}
\end{figure}

\subsection{Help and exit}
\label{subsec:shell-help-and-exit}

\texttt{help} without arguments prints the main command list, split into
ordinary commands and the grouped commands that have topic help. The grouped
topics currently include \texttt{sim}, \texttt{config}, \texttt{plugin},
\texttt{model}, \texttt{trace}, \texttt{execute-script}, \texttt{load},
\texttt{save}, \texttt{facade}, \texttt{licence}, \texttt{data},
\texttt{component}, and \texttt{experiment}.

\texttt{help <topic>} expands the selected topic and shows the current
subcommand vocabulary. This is the safest way to discover the command surface
without reading the source.

\texttt{exit} sets the exit flag and returns control to the outer loop. The
current implementation prints a quit message and stops the interactive loop.

\subsection{Metadata, identity, and trace}
\label{subsec:shell-metadata}

\texttt{facade} reports the simulator name and version metadata through
\texttt{SimulatorFacade}. The subcommands \texttt{get-version},
\texttt{get-version-number}, and \texttt{get-name} are confirmed in the code.

\texttt{info} exposes the current model metadata. The supported actions include
reading and updating the model name, analyst name, description, project
title, version, and changed flag.

\texttt{licence} prints licence details and the current model or runtime
limits. The shell supports show, limits, activation-code, the activation-code
search and insertion/removal actions, and the model/entity/host/thread limit
queries.

\texttt{trace} reads or changes the current trace level. The shell accepts
\texttt{show} and \texttt{level <1-9>} and stores the chosen level in the
underlying trace manager.

\subsection{Plugin management}
\label{subsec:shell-plugins}

\texttt{plugin} inspects installed plugins, prints their count, shows detailed
information for a selected plugin, autoloads a plugin list file, displays a
language template, validates a candidate library, checks system dependencies,
discovers plugin filenames, and navigates the plugin list.

The current shell has a safety gate around dependency installation. When a
plugin load path requires system packages, the shell only offers automatic
installation when every missing dependency can be installed automatically and
stdin is interactive. In non-interactive mode, it refuses to run the install
command automatically and asks the user to retry manually.

\subsection{Model and simulation}
\label{subsec:shell-model-simulation}

The model commands are split between grouped and convenience forms.
\texttt{model} supports creating, removing, checking, showing, loading, saving,
and navigating the current model list. The top-level \texttt{load} and
\texttt{save} commands are convenience aliases for the current model file
operations.

\texttt{sim} is the main simulation inspection and control command. It can
show the current state, start, pause, step, stop, read counters and timing
values, change the replication configuration, toggle pause behavior, and read
or update reporting options.

\texttt{config} is the configuration-oriented command family. It focuses on
replication count, replication length, warm-up, base time unit, terminating
condition, pause behavior, and the report toggles.

\texttt{simul} remains as a compact control alias in the current code. The
shell also keeps the lower-level report and breakpoint helpers exposed by
\texttt{sim} for users who need fine-grained control during debugging.

\begin{figure}[htbp]
\centering
% FIGURE-SPEC-BEGIN
% ID: fig:shell-model-loaded
% Source of truth:
%   - source/applications/shell/GenesysShell.cpp
%   - models/AirportSecurityExample.gen
%   - source/kernel/simulator/SimulatorFacade.h
% Purpose:
%   Show the shell after a real model has been loaded and inspected.
% Required visual content:
%   - the model load command
%   - a confirmation that the model loaded
%   - a follow-up inspection command such as model show, model check, or sim show
% Accessibility description:
%   A terminal screenshot that proves the shell can move from startup to a loaded model.
% Validation:
%   The observed model path and output must match a repository model file.
% FIGURE-SPEC-END
\figureplaceholder{figs/generated/shell-model-loaded}
\caption{Shell session after loading a model file.}
\label{fig:shell-model-loaded}
\end{figure}

Figure~\ref{fig:shell-model-loaded} should be read as the minimum proof that
the shell can move from startup into a loaded-model workflow.

\subsection{Inspection helpers}
\label{subsec:shell-inspection}

\texttt{data} inspects the current model data definitions. The current code can
count definitions, list them by class name, retrieve individual entries, clear
the set, and read or update the changed flag.

\texttt{component} inspects the current model components. It can count, show,
list, clear, search by id or name, move through the list, and update the
changed flag.

\texttt{experiment} inspects simulation experiments. It can count, show the
current experiment, create a new one, save or load an experiment file, and
navigate the experiment list.

\section{Loading, validation, and execution}
\label{sec:shell-loading-validation-execution}

There are two model-loading paths in the current shell:
\begin{itemize}
\item \texttt{model load <filename>};
\item \texttt{load <filename>}.
\end{itemize}

Both paths call the current facade loader and print either \texttt{Model loaded}
or an error message. The file path is used as provided, so the command should be
run from a directory where that path resolves correctly.

\texttt{model check} validates the loaded model through the facade. The shell
prints a notice if the check reports errors. This is the current validation
step for the terminal workflow; it confirms that the loaded model is accepted by
the code, not that the scenario is scientifically sufficient.

\texttt{sim start}, \texttt{sim pause}, \texttt{sim step}, and \texttt{sim stop}
control execution. The \texttt{sim show} command prints the current simulation
state, including the configured parameters and runtime flags.

\begin{figure}[htbp]
\centering
% FIGURE-SPEC-BEGIN
% ID: fig:shell-terminal-result
% Source of truth:
%   - source/applications/shell/GenesysShell.cpp
%   - source/kernel/simulator/SimulationReporterDefaultImpl1.cpp
%   - source/kernel/simulator/model/ModelSimulation.cpp
% Purpose:
%   Show a real terminal result produced after loading and checking a model, or after starting a short run.
% Required visual content:
%   - the command that was executed
%   - the immediate terminal response
%   - any follow-up line that proves the command had effect
% Accessibility description:
%   A result screenshot that proves the shell is doing real work, not only echoing prompts.
% Validation:
%   The shown output must be produced by the current shell or facade code.
% FIGURE-SPEC-END
\figureplaceholder{figs/generated/shell-terminal-result}
\caption{Terminal result after a real shell command.}
\label{fig:shell-terminal-result}
\end{figure}

The terminal result in Figure~\ref{fig:shell-terminal-result} should be read
as evidence of the current command path. It must come from the real shell
output, not from an edited mockup.

\section{Output, return, and automation}
\label{sec:shell-output-automation}

The shell prints its own diagnostics directly to standard output. There is no
separate structured return object for command success or failure, so the user
must read the text output to understand what happened.

The current process exit path is conservative: \texttt{main} returns \texttt{0}
after the shell loop ends. In practice, this means that command failures are
communicated through text rather than by a rich exit-status protocol.

Automation is supported in three ways:
\begin{itemize}
\item additional command strings passed on the argv list;
\item \texttt{execute-script <filename>} for line-oriented shell scripts;
\item \texttt{bash <command>} for raw host-shell execution.
\end{itemize}

The argv path executes each argument string as a command string in order. This
is useful for compact command sequences, but it is not a full shell parser.
When a command line needs spaces, quote the entire string at the outer shell
layer or prefer \texttt{execute-script}.

\texttt{execute-script} ignores blank lines and comment lines that begin with
\texttt{\#}. Each executed line is echoed with the \texttt{\$execute-script>}
prefix before it is routed through the same command parser used by the
interactive shell.

\texttt{bash} forwards the collected argument text to \texttt{system()}. That
makes it powerful but also unsafe for untrusted input. It is intended for local
automation, not for a security boundary.

\begin{figure}[htbp]
\centering
% FIGURE-SPEC-BEGIN
% ID: fig:shell-automation
% Source of truth:
%   - source/applications/shell/GenesysShell.cpp
%   - source/applications/shell/main.cpp
% Purpose:
%   Show a small scripted or argv-driven session that demonstrates automation.
% Required visual content:
%   - an argv-driven command or execute-script run
%   - the echoed command lines
%   - a real output line from the shell
% Accessibility description:
%   A screenshot that demonstrates the non-interactive shell path.
% Validation:
%   The session must use a command path confirmed in the current source.
% FIGURE-SPEC-END
\figureplaceholder{figs/generated/shell-automation}
\caption{Automated shell execution through argv or script input.}
\label{fig:shell-automation}
\end{figure}

Figure~\ref{fig:shell-automation} is the right place to document the scripted
path because it proves the shell can be used for lightweight automation without
opening the GUI.

\section{Current limits}
\label{sec:shell-limits}

The shell is intentionally narrow and has several limits that should be stated
explicitly.
\begin{itemize}
\item Command parsing is whitespace based; there is no general quoted-argument
parser inside the shell itself.
\item Many commands expect exactly one path token, so file names with spaces are
not a supported pattern.
\item Failures are mostly textual and do not map to a detailed exit-status
protocol.
\item \texttt{bash} executes host commands directly and should only be used on
trusted local input.
\item \texttt{plugin} autoload can prompt for dependency installation only when
the input is interactive and all missing dependencies are eligible for
automatic installation.
\item The shell is a command surface over the current simulator code; it does
not add a separate policy layer for security, scientific validation, or release
approval.
\end{itemize}

These limits are deliberate. They keep the shell small, predictable, and easy
to audit, while leaving stronger policy decisions to the surrounding
documentation and governance process.
